% Folder 156-8, batch 02 — archivists' pages 21–40 (batch PDF pages 1–20).
% Pass: Opus 5.5 (claude-opus-5-5), 2026-10-02 — first pass, unchecked against the pages by a human.
%
% Page numbering follows the batch mapping (batch page 1 = archivists' page 21).
% The pencilled archivists' numbers visible on the scans run one ahead of it
% (the sheet given here as page 35 is pencilled « 36 », page 40 « 41 »);
% to be reconciled with the folder's other batches. Grothendieck's own
% pagination (pp. 13–32) is recorded in a note at the head of each page.
\documentclass[11pt,a4paper]{article}
\input{../preamble/grothendieck}

\folder{156-8}
\batch{2}
\pages{21}{40}
\foldertitle{[Chapitre] VIII. Analysis situs (quatrième mouture) : notes manuscrites (26/06-04/07/1986).}
\dating{1986}
\watermark{Édition de démonstration}

\begin{document}

\page{21}
\note{p.~13 de l'auteur.}
Prise de vue « synthétique ». Structure sur deux ensembles de base
$\mathcal{M}$ (multistructures), $\mathfrak{F}$ (figures).

Structure
\begin{itemize}
  \item[a)] Relation d'incidence $X \lhd F$ entre $\mathcal{M}$, $\mathfrak{F}$
    (\uncertain{suffisante} à définir $\leq$ sur $\mathcal{M}$ et sur $\mathfrak{F}$)
  \item[b)] Relations $\ll$, $|\circ|$ sur $\mathcal{M}$
\end{itemize}

Pour tout $F \in \mathfrak{F}$, on pose
\[
  \widetilde{F} = \{ X \in \mathcal{M} \mid X \lhd F \}
  \qquad \text{On pose } F \leq G \overset{\mathrm{def}}{\Longleftrightarrow}
  \widetilde{F} \subset \widetilde{G}
\]

\struck{\emph{At M 1} $\widetilde{F} = \widetilde{G} \Rightarrow F = G$, i.e.
$\mathfrak{F} \to \mathfrak{P}(\mathcal{M})$, $F \mapsto \widetilde{F}$ est
une injection, ou encore la relation $\leq$ sur $\mathfrak{F}$ est un
\emph{ordre}.}

\struck{\emph{At M 2} $\widetilde{\mathfrak{F}}$ \ill{} est une \ill{}
\ill{} qui \uncertain{connaisse} $\mathcal{M}$ et \uncertain{sépare} les points
de $\mathcal{M}$, i.e. a) $\forall X \in \mathcal{M}$,
$\mathfrak{F}^{X} = \{ F \in \mathfrak{F} \mid X \lhd F \}$ a un plus petit
élément (on l'appelle la \ill{} \uncertain{multistructure} $X$)}
\note{tout ce passage est barré de deux longs traits obliques ; on y lit encore, au-dessus de la ligne de At M 2, une variante elle-même biffée : « ens. de parties de $\mathcal{M}$ quasi-admissible ».}
\marginal{Ainsi,}

\emph{At M 1} \add{b)} La relation $X \lhd F$ est \uncertain{séparante} sur
$\mathcal{M}$ et sur $\mathfrak{F}$, i.e.
\begin{itemize}
  \item[a)] L'application $F \longmapsto \widetilde{F}$,
    $\mathfrak{F} \longrightarrow \mathfrak{P}(\mathcal{M})$ est injective
    [i.e. encore, $\leq$ est une relation d'\emph{ordre}]
  \item[b)] Si $X, Y \in \mathcal{M}$ tels que $X \neq Y$, alors
    $\exists F \in \mathfrak{F}$ avec $X \in \widetilde{F}$,
    $Y \notin \widetilde{F}$ ou l'inverse (i.e.
    $\widetilde{\mathfrak{F}} \subset \mathfrak{P}(\mathcal{M})$ sépare les
    points de $\mathcal{M}$)
\end{itemize}

\emph{At M 2} La famille de parties $\widetilde{\mathfrak{F}}$ dans
$\mathcal{M}$ est quasi-admissible, i.e. $\forall X \in \mathcal{M}$,
\struck{\ill{}} posant
$\mathfrak{F}^{X} = \{ F \in \mathfrak{F} \mid X \lhd F \}$, $\mathfrak{F}^{X}$
a un plus petit élément (appelé la \emph{figure} \emph{élémentaire} associée à
la multistructure $X$)

Ainsi \struck{$X \lhd F$ équivaut à}
\[
  X \lhd F \Longleftrightarrow F_{X} \leq F \qquad
  (\Longleftrightarrow X \in \widetilde{F})
\]
de sorte que la condition At M 1 b), compte tenu de At M 1 a), équivaut à
\[
  \mathcal{M} \longrightarrow \mathfrak{F} \qquad X \longmapsto F_{X}
  \quad \text{est injectif.}
\]

\emph{At M 3}
\begin{itemize}
  \item[a)] Soit \struck{\ill{}} $\mathfrak{L}_{i}, \mathfrak{L} \in
    \widetilde{\mathfrak{F}}$, avec $\mathfrak{L}_{i} \subset \mathfrak{L}$.
    Alors $\bigcup \mathfrak{L}_{i} \in \widetilde{\mathfrak{F}}$ (stabilité de
    $\widetilde{\mathfrak{F}}$ par $\bigcup$ \uncertain{majorées})
  \item[b)] $\mathfrak{F} \neq \emptyset$
    ($\Longleftrightarrow \emptyset \in \widetilde{\mathfrak{F}}$
    $\Longleftrightarrow \widetilde{\mathfrak{F}} \neq \emptyset$
    $\Longleftrightarrow \mathfrak{F}$ a un plus petit élément)
\end{itemize}
\note{la lettre notée $\mathfrak{L}$ est une majuscule cursive bouclée, soulignée ; lecture incertaine.}
\marginal{résulte du reste si $\mathcal{M} \neq \emptyset$}

\page{22}
\note{p.~14 de l'auteur.}
Les conditions At M 1 -- At M 3 \add{plus At M 4 plus bas} ne concernent que
la relation $X \lhd F$. Allons un peu plus loin.
\marginal{$\lhd$ $\leftarrow$}

Si $X, Y \in \mathcal{M}$, on a tautologiquement
\begin{align*}
  F_{X} \leq F_{Y}
  &\Longleftrightarrow \forall\, \text{\struck{$\mathfrak{L} \in \mathfrak{F}$}}\
    F \in \mathfrak{F},\ Y \lhd F \Rightarrow X \lhd F \\
  &\Longleftrightarrow \forall\, \mathfrak{L} \in \widetilde{\mathfrak{F}},\
    Y \in \mathfrak{L} \Rightarrow X \in \mathfrak{L} \\
  &\Longleftrightarrow X \lhd F_{Y} \Longleftrightarrow X \in \widetilde{F}_{Y}
\end{align*}
On écrit alors
\[
  X \leq Y ,
\]
c'est une relation de préordre, \uncertain{induite} par la relation de
préordre associée à la famille $\widetilde{\mathfrak{F}}$ de parties dans
$\mathcal{M}$. Par At M 1 b) c'est une relation d'ordre. Ceci dit, on a aussi
\marginal{On écrit $\widetilde{X} = \widetilde{F}_{X} = \{ Y \mid Y \lhd F_{X} \}
= \mathcal{M}_{\leq X}$}

\struck{\emph{Prop.} Si $F \in \mathfrak{F}$, $\widetilde{F}$ est une partie
fermée de $\mathcal{M}, \leq$.} \struck{Si \ill{} est une partie fermée de
$\mathfrak{L}$, alors $X \in \widetilde{F}$}
\note{deux lignes barrées ; au-dessus de la première est écrit $\mathfrak{L} = \widetilde{F}$.}

\emph{Prop.} Soit $\mathfrak{L} \in \widetilde{\mathfrak{F}}$, alors toute
partie fermée \struck{de} \add{\uncertain{dans}} $\mathfrak{L}$ est dans
$\widetilde{\mathfrak{F}}$.

Ceci \struck{\ill{}} \struck{At} est formellement plus fort que At M 3 a), et
lui est même équivalent (compte tenu du reste). Si j'ai laissé At M 3 a) sous
sa forme \uncertain{apparemment} plus faible, c'est pour \uncertain{l'exprimer}
\ill{} en termes de $(\mathfrak{F}, \leq)$, comme une condition d'existence de
Sup majorés.

\emph{At M 4} \struck{Soit $F, G \in \mathfrak{F}$. Pour que}
Soit $F, G \in \mathfrak{F}$, on pose
\begin{align*}
  F \lessgtr G &\overset{\mathrm{def}}{\Longleftrightarrow}
    \{F, G\} \text{ majoré dans } \mathfrak{F} \Longleftrightarrow
    F \cup G \text{ \uncertain{existe}} \\
  &\Longleftrightarrow \widetilde{F} \cup \widetilde{G} \in \widetilde{\mathfrak{F}} \\
  X \lessgtr Y &\overset{\mathrm{def}}{\Longleftrightarrow} F_{X} \lessgtr F_{Y}
    \text{ i.e. } \widetilde{X} \cup \widetilde{Y} \in \widetilde{\mathfrak{F}}
    \quad [\widetilde{X} = \mathcal{M}_{\leq X}]
\end{align*}

\emph{At M 4}
\[
  F \lessgtr G \Longleftrightarrow \forall X \in \widetilde{F},\
  Y \in \widetilde{G}, \text{ on a } X \lessgtr Y
\]

\struck{\ill{}} Pour les relations \ill{} $\ll$, $|\circ|$ sur $\mathcal{M}$,
on pose

\emph{At M 5}
\[
  X \lessgtr Y \Longleftrightarrow \forall X' \in \widetilde{X} \setminus
  \widetilde{X} \cap \widetilde{Y},\
  Y' \in \widetilde{Y} \setminus \widetilde{X} \cap \widetilde{Y},
\]
on a $X' \mathrel{|\circ|} Y'$.
\note{le second membre de At M 5 est écrit à l'étroit en bas de page ; l'ensemble où court $Y'$ se lit mal, la lecture $\widetilde{Y} \setminus \widetilde{X} \cap \widetilde{Y}$ suit celle de la p.~26.}

\page{23}
\note{p.~15 de l'auteur.}
\emph{Plus} les axiomes \add{des magasins} (reliant $\leq$, $\ll$, $\circ$ sur
$\mathcal{M}$), savoir Mag 1 -- Mag 4.

Ainsi, \uncertain{qu'on parte} de $(\mathfrak{F}, \leq, \ll, |\circ|)$,
\uncertain{ou} de $(\mathfrak{F}, \mathcal{M}, \lhd, \ll, |\circ|)$, il nous
faut, en plus de Mag 1 -- Mag 4, \uncertain{cinq} axiomes. Si on part de
$(\mathcal{M}, \leq, \ll, |\circ|, \mathfrak{F})$, il nous faut un seul axiome
(\uncertain{itinéraires}…) supplémentaire, concernant le choix de
$\mathfrak{F}$.

Tout ceci concerne les ateliers que je devrais \add{\uncertain{plutôt}} appeler
« spécieux ». Il faut maintenant ouvrir une parenthèse sur les « ateliers
“polyidéaux” », qui représentent une \struck{structure} \add{variante} très
légèrement différente.

\emph{Définition.} Le magasin $(\mathcal{M}, \leq$ \ill{}$)$ est dit
\add{pré}\emph{polyidéal} \struck{si $\forall X \in \mathcal{M}$,
$\widetilde{X}$ est un ens. ordonné \ill{} polyidéal, i.e.} si deux éléments
$X, Y$ de $\mathcal{M}$ qui sont majorés, ont un Inf \add{dans $\mathcal{M}$}.
\struck{\ill{} est dit polyidéal si de plus \ill{} sont majorés, \ill{}
existe}
\note{la définition est marquée d'un trait vertical en marge ; la première rédaction, raturée, s'enchevêtre avec la seconde, et seule celle-ci se lit en entier.}

\struck{\emph{NB} Si les $\mathcal{M}_{a \leq \cdot \leq b}$ \ill{} est contenu
dans les \ill{} plus grand, si $\mathcal{M}$ est \uncertain{polyidéal}
\ill{} et si les $\mathcal{M}_{a \leq \cdot \leq b}$ sont inductifs, alors
deux \ill{}}
\note{passage encadré et barré de traits obliques.}

On introduit \struck{\ill{}} les relations plus fortes que $\ll$ (pour les
magasins, pas même prépolyidéaux…)
\[
  X \ll_{\mathrm{pol}} Y \Longleftrightarrow X \ll Y, \text{ et }
  \forall Y' \in \widetilde{Y},\ X_{Y'} \overset{\mathrm{def}}{=}
  \widetilde{X} \cap \mathrm{Ombr}(Y')
\]
est vide ou a un plus grand élément (i.e. $X_{Y'}$ est une figure élémentaire)

\marginal{En fait $X_{Y'} = X \cdot Y' = \mathrm{Inf}^{\ll}(X, Y')$ dans
$\mathfrak{F}_{\mathcal{M}}$, dans le formalisme des figures…}
\note{remarque encadrée en bas à gauche de la page.}

\page{24}
\note{p.~16 de l'auteur.}
La condition que $\mathcal{M}$ soit \add{pré}polyidéal \struck{\ill{}} équivaut à
dire que
\[
  X \leq Y \Longrightarrow X \ll_{\mathrm{pol}} Y
\]
\note{le signe de droite est un $\ll$ tracé d'un trait plus lourd, peut-être repassé ; la lecture $\ll_{\mathrm{pol}}$ est celle de l'indice « pol » écrit dessous.}
ce qui est nécessaire pour que
\[
  (\mathcal{M}, \leq, \ll_{\mathrm{pol}}, |\circ|)
\]
soit encore un magasin. Est-ce suffisant ?

1°) \emph{Mag 1.} Il reste à vérifier que $\ll_{\mathrm{pol}}$ est une relation
d'ordre, et tout revient à la transitivité
\[
  X \ll_{\mathrm{pol}} Y \ll_{\mathrm{pol}} Z \Longrightarrow X \ll_{\mathrm{pol}} Z .
\]
Soit $Z' \leq Z$, alors que $Y \ll_{\mathrm{pol}} Z$ \struck{\ill{}} la
\uncertain{sous-figure} induite $Y^{*}_{Z'}$ est vide, ou de la forme $Y'$. Dans
le premier cas, la figure induite par $Y_{Z'}$ sur $X$ est vide, OK, dans le
deuxième elle est vide ou de la forme $X'$, par $\ll_{\mathrm{pol}}$, OK aussi.
(On utilise \uncertain{transitivité} \ill{} \ill{})

2°) \emph{Mag 2.} Soit $X \ll_{\mathrm{pol}} Y$, prouvons que l'ens.
\[
  \widetilde{Y}^{X,\mathrm{pol}} = \{ Z \in \widetilde{Y} \mid X \ll_{\mathrm{pol}} Z \}
\]
a un plus petit élément, pour ceci, il suffit de noter que l'inclusion
évidente $\subset$ est une égalité
\[
  \widetilde{Y}^{X,\mathrm{pol}} = \widetilde{Y}^{X} ,
\]
\struck{et on fait une égalité}, i.e. que \struck{$X \leq$}
\[
  X \ll Y' \leq Y \Longrightarrow X \ll_{\mathrm{pol}} Y' ,
\]
c'est évident par définition (vu que les $Y'' \leq Y'$ sont aussi
$\in \widetilde{Y}$).

\struck{\ill{}} \emph{NB} On a constaté que
\[
  X \mathrel{\mathring{\ll}_{\mathrm{pol}}} Y \Longleftrightarrow
  X \ll_{\mathrm{pol}} Y \text{ et } X \mathrel{\mathring{\ll}} Y .
\]

\page{25}
\note{p.~17 de l'auteur.}
3°) \emph{Mag 3} \ill{} trivial, puisque $|\circ|$ \uncertain{inchangé} par
\uncertain{changement} ; et $\mathring{\ll}_{\mathrm{pol}} \Rightarrow \mathring{\leq}$.
\note{le signe à droite de l'implication est un $\leq$ surmonté d'un petit rond, avec un trait raturé sous lui ; lecture incertaine.}

4°) \emph{Mag 4} \struck{aussi} \uncertain{idem} par \uncertain{changement} (on ne
fait intervenir que $\leq$ et $|\circ|$)

\emph{Définition.} Un magasin est dit \struck{\ill{}} \emph{polyidéal}, si
$X \ll Y \Rightarrow X \ll_{\mathrm{pol}} Y$ (auquel cas il est \struck{poly}
aussi prépolyidéal, comme on voit en prenant $X \leq Y$),
\struck{Si $\mathcal{M}, \leq, \ll, |\circ|$} donc si $\ll$ et
$\ll_{\mathrm{pol}}$ sont égales dans $\mathcal{M}$. Le magasin
\[
  \mathcal{M}^{\mathrm{pol}} = (\mathcal{M}, \leq, \ll_{\mathrm{pol}}, |\circ|)
\]
est dit le magasin polyidéal associé au magasin prépolyidéal $\mathcal{M}$
\add{(il est bien polyidéal)}. On a \struck{\ill{}
$(\mathcal{M}^{\mathrm{pol}})^{\mathrm{pol}} = \mathcal{M}^{\mathrm{pol}}$, et}
$\mathcal{M} = \mathcal{M}^{\mathrm{pol}}$ ssi $\mathcal{M}$ est polyidéal.

\struck{Toujours} Mais si $\mathcal{M}$ est déjà polyidéal, i.e. il n'y a pas
à faire de différence entre $\ll$ et $\ll_{\mathrm{pol}}$, il
\struck{faut} \uncertain{faut} distinguer deux notions de compatibilité,
$\lessgtr$ et $\lessgtr_{\mathrm{pol}}$.
\note{le signe de compatibilité, un $<$ et un $>$ croisés posés sur un trait, est rendu ici et dans toute la suite par $\lessgtr$.}

\emph{Définition.} Soit $\mathcal{M}$ magasin. On pose
\[
  X \lessgtr_{\mathrm{pol}} Y \overset{\mathrm{def}}{\Longleftrightarrow}
  X \lessgtr Y, \text{ et } \forall X' \in \widetilde{X},\ Y' \in \widetilde{Y},
\]
$\widetilde{X}' \cap \widetilde{Y}'$ est vide ou a un plus grand élément.

Donc $\mathcal{M}$ est prépolyidéal ssi pour tout $X \in \mathcal{M}$, on a
\struck{$X$ et $Y, Z \leq X$, on a $Y \lessgtr_{\mathrm{pol}} Z$. Donc ce cas}
$X \lessgtr_{\mathrm{pol}} X$, donc $Y \lessgtr_{\mathrm{pol}} Z$ pour
$Y, Z \leq X$.

\marginal{\ill{} \ill{} $X \lessgtr_{\mathrm{pol}}$ \ill{} $Y \lessgtr$ \ill{}}
\note{note marginale tournée, embrassée d'une accolade qui la rattache aux deux lignes précédentes ; seuls quelques signes s'en lisent.}

On introduit \uncertain{donc} \struck{Fig}
\[
  \mathrm{Figpol}(\mathcal{M}) \subset \mathrm{Fig}(\mathcal{M})
\]
(\uncertain{anciennement} $\mathcal{F}_{\mathcal{M}}$) formée des figures $F$ telles que
$\forall X, Y \in \widetilde{F}$, on a $X \lessgtr_{\mathrm{pol}} Y$,

\page{26}
\note{p.~18 de l'auteur.}
i.e. des \struck{figures} parties $F \subset \mathcal{M}$ telles que
\begin{itemize}
  \item[a)] $F$ fermé par $\leq$
  \item[b)] $X, Y \in F \Longrightarrow X \lessgtr_{\mathrm{pol}} Y$, i.e.
\end{itemize}
\[
  \left\{
  \begin{array}{l}
    1^{\circ})\ \forall X' \in \widetilde{X} \setminus \widetilde{X} \cap \widetilde{Y},\
      Y' \in \widetilde{Y} \setminus \widetilde{X} \cap \widetilde{Y}
      \Longrightarrow X' \mathrel{|\circ|} Y' \\
    2^{\circ})\ \forall X' \in \widetilde{X},\ Y' \in \widetilde{Y},\
      \widetilde{X}' \cap \widetilde{Y}' \text{ est vide ou a un plus grand élément.}
  \end{array}
  \right.
\]
\note{dans 1°, le second ensemble a d'abord été écrit avec $\widetilde{X}$, raturé et récrit $\widetilde{Y}$.}

\struck{Ceci posé, on trouve}

\emph{NB} Même si $\mathcal{M}$ est polyidéal, on a en général une inclusion
\emph{stricte} $\mathrm{Figpol}(\mathcal{M}) \subsetneq \mathrm{Fig}(\mathcal{M})$.

\drawing{un petit cercle portant deux points marqués, l'un en haut, l'autre en bas.}
Cette figure circulaire est dans $\mathrm{Fig}(\mathcal{M})$, mais non dans
$\mathrm{Figpol}(\mathcal{M})$, pour $\mathcal{M}$ magasin polyidéal convenable.

\drawing{dans la marge gauche, deux triangles : dans le premier un arc de cercle s'appuie sur la base, dans le second un arc joint deux sommets de la base.}
\marginal{Exemples analogues avec $X \ll Y$ mais non $X \ll_{\mathrm{pol}} Y$
(2 dimensions)}

\emph{Scholie.} Les propriétés \struck{formelles} \add{algébriques} de
$\mathcal{F} = \mathrm{Figpol}(\mathcal{M})$, muni de $\leq$,
$\ll_{\mathrm{pol}}$, $\lessgtr_{\mathrm{pol}}$, et ses relations avec
$\mathcal{M}$ (\uncertain{via} des relations induites) seront essentiellement
les mêmes, que pour $\leq$, $\ll$, $\lessgtr$ \struck{pol, pol}. La seule
différence est dans l'explicitation de $\lessgtr_{\mathrm{pol}}$
\uncertain{vis-à-vis} $\lessgtr$, en termes de la relation de disjonction
intérieure $|\circ|$.

On devrait prendre en revue \struck{toutes} les propositions pertinentes…
\note{ce paragraphe est marqué d'un trait vertical en marge.}

Je signale seulement celle-ci

\emph{Prop.} Si $F, G \in \mathrm{Figpol}$, alors \struck{$\{F, G\}$ majoré
\ill{}} conditions équivalentes
\begin{itemize}
  \item[(i)] $\{F, G\}$ majoré par $\leq$
  \item[(ii)] $\mathrm{Sup}^{\leq}(F, G)$ existe
  \item[(ii')] $F \cup G$ (réunion dans $\mathfrak{P}(\mathcal{M})$) $\in \mathrm{Figpol}(\mathcal{M})$
  \item[(iii)] $\forall X \in \widetilde{F}$, $Y \in \widetilde{G}$, on a $X \lessgtr_{\mathrm{pol}} Y$.
\end{itemize}
\marginal{On écrit alors $F \lessgtr_{\mathrm{pol}} G$}
\note{la remarque est écrite à droite des quatre conditions, reliée par une accolade.}

Dans certains contextes, on laissera tomber l'indice pol à $\lessgtr$, $\ll$,
quand il est bien entendu qu'il ne s'agit que des relations « polyidéales ».

\emph{Atelier polyidéal.} Se définit essentiellement \add{(en termes d'un
magasin polyidéal et de $\mathrm{Figpol}(\mathcal{M})$)}

\page{27}
\note{p.~19 de l'auteur.}
de la même façon qu'un atelier « spécieux » (ou ordinaire) se définit en termes
d'un magasin quelc. et de $\mathrm{Fig}(\mathcal{M})$. En termes d'un
\[
  \mathfrak{F} \subset \mathfrak{P}(\mathcal{M}), \text{ ou plutôt }
  \mathfrak{F} \subset \mathrm{Figpol}(\mathcal{M}),
\]
les axiomes pertinents, qui paraphrasent Mag at (p.~10) \uncertain{sont}

\emph{Mag at pol}
\begin{itemize}
  \item[a)] $F \in \mathfrak{F}$, $G \leq F$ (i.e. $G$ partie fermée de $F$ pour
    $\leq$ dans $\mathcal{M}$) $\Longrightarrow G \in \mathfrak{F}$
  \item[b)] $F, G \in \mathfrak{F}$, $F \lessgtr_{\mathrm{pol}} G \Longrightarrow F \cup G \in \mathfrak{F}$
  \item[c)] $\forall X \in \mathcal{M}$, $\widetilde{X} \in \mathfrak{F}$
  \item[c')] $\mathfrak{F} \neq \emptyset$ ($\Longleftrightarrow \emptyset \in \mathfrak{F}$).
\end{itemize}

\emph{NB} Un « atelier spécieux » \struck{est} \add{aussi} un « atelier
polyidéal » ssi le magasin associé est polyidéal \struck{\ill{}} (i.e.
$X \ll Y \Rightarrow X \ll_{\mathrm{pol}} Y$) \struck{\ill{}} (chose \add{qui}
\uncertain{ici va de soi}), et si \emph{de plus},
\[
  X \lessgtr Y \Longrightarrow X \lessgtr_{\mathrm{pol}} Y
\]
(ce qui est beaucoup plus exceptionnel(?)). Je suis \uncertain{amené} à faire
\ill{} \uncertain{remarques} que ce sont des notions très voisines, mais
disjointes, et que
\begin{itemize}
  \item[a)] Les notions « spécieuses » \struck{\ill{}} $\ll$, $\lessgtr$ sont
    plus commodes quand il s'agit de travailler avec des figures aussi
    générales que possibles,
\end{itemize}
et que
\begin{itemize}
  \item[b)] Les notions « polyidéales » $\ll_{\mathrm{pol}}$,
    $\lessgtr_{\mathrm{pol}}$ sont plus commodes quand il s'agit de travailler
    avec des figures, \uncertain{par exemple} des « découpages » \add{d'autres}
    en « morceaux » aussi « standard » que possible, avec des « ajustements »
    \struck{\ill{}} des « morceaux » aussi \struck{A} standard que possible.
\end{itemize}

\marginal{Cela provient \uncertain{seulement} du fait qu'on a pris $\ll$ comme
relation primitive \uncertain{plutôt} \ill{} mais non $\lessgtr$. Il y a un
\uncertain{certain} arbitraire…}
\note{note marginale écrite en oblique dans l'angle inférieur gauche, rattachée au passage sur $X \lessgtr Y \Rightarrow X \lessgtr_{\mathrm{pol}} Y$.}

\page{28}
\note{p.~20 de l'auteur.}
Dans le magasin, la différence des deux points de vue se réduit au minimum :
On ne la voit pas dans la notion des données, \struck{et dans les axiomes},
à cela près que dans le cas « polyidéal », on demande des axiomes plus
stricts, savoir, en plus de Mag 1 -- Mag 4, l'axiome supplémentaire

\emph{Mag pol} Pour $X \ll Y$ et $Y' \in \widetilde{Y}$,
$\widetilde{X} \cap \mathrm{Omb}(Y')$ est vide, ou a un plus grand élément,

qui exprime que $X \ll Y \Longrightarrow X \ll_{\mathrm{pol}} Y$.

Cela n'empêche qu'il faudra, par la suite, en plus de $\lessgtr$, (ou
\uncertain{au lieu} de $\lessgtr$), introduire la relation
$\lessgtr_{\mathrm{pol}}$ \ill{}

\note{un long trait oblique sépare ce qui précède de ce qui suit.}

Je reviens encore sur l'axiomatique \uncertain{basée} en termes de
Mag 1 : Mag 4, en notant que la connaissance de $\mathring{\ll}$, $|\circ|$
redonne celle de $\ll$ par l'équivalence
\[
  (*) \qquad X \ll Y \Longleftrightarrow \exists Y',\
  X \mathrel{\mathring{\ll}} Y' \leq Y .
\]
On se demande, dès lors, quand \struck{At} \add{\uncertain{une}}
\struck{d'une structure donnée} \add{rel. d'ordre}
\[
  (\mathcal{M}, \leq, \mathring{\ll}, |\circ|) ,
\]
\struck{\ill{}} celle-ci provient d'un magasin. Il faut bien examiner les
conditions pour que
\note{la phrase se poursuit à la page suivante.}

\page{29}
\note{p.~21 de l'auteur.}
\begin{itemize}
  \item[a)] $(\mathcal{M}, \leq, \ll, |\circ|)$, avec $\ll$ défini par $(*)$,
    soit un magasin, et
  \item[b)] La relation de \uncertain{recouvrement} immédiat dans ce magasin
    est la relation de départ $\mathring{\ll}$.
\end{itemize}

Il faut supposer déjà
\note{les numéros (1), (2), (3) et (3') de cette page sont cerclés dans le manuscrit, de même que le (d) de la p.~35.}

(1) $\leq$, $\mathring{\ll}$ \struck{\ill{}} \add{\uncertain{relations}}
relations d'ordre,

et pour vérifier Mag 1, il reste à vérifier que $X \ll Y$ est une relation
d'ordre. Ceci se ramène à la situation (qu'il faut \uncertain{sans doute}
prendre comme axiome)
\[
  (2) \qquad X \leq Y \mathrel{\mathring{\ll}} Z \Longrightarrow
  \text{\struck{$\exists$}}\ \exists Z' \text{ avec } X \mathrel{\mathring{\ll}} Z' \leq Z ,
\]
avec \struck{\ill{}} l'\emph{unicité} en $Z'$. \struck{Dans le premier cas}
\note{une variante interlinéaire, sous « avec l'unicité », se lit « unicité si \ill{} \ill{} ».}

\emph{\struck{Mag 2}} (\struck{2}\add{3}') si $X \leq Y$ telles qu'il
existe $Y'$ avec $X \mathrel{\mathring{\ll}} Y' \leq Y$, alors $Y'$ est unique.
\marginal{\uncertain{inutile}}

Sous ces conditions, Mag 2 devient automatique, ainsi que le fait que si
$X \mathrel{\mathring{\ll}} Y$, alors \struck{$Y$} \struck{\ill{}}
$X \ll Y' \leq Y \Rightarrow Y' = Y$ (i.e. \struck{$X$} est \uncertain{bien
défini} immédiatement en $Y$, au sens des \uncertain{rel.}\ $\leq$, $\ll$)

\emph{Mag 3} reste une condition, (3). Mag 4 \uncertain{aussi}.

En résumé, l'axiomatique devient

\emph{Mag 1'} $X \leq Y$ et $X \mathrel{\mathring{\ll}} Y$ sont des relations
d'ordre.

\emph{Mag 2'}
\begin{itemize}
  \item[a)] Soient $X, Y$ tels qu'il existe $Y'$, avec
    $X \mathrel{\mathring{\ll}} Y' \leq Y$. Alors cet $Y'$ est uniquement
    déterminé.
  \item[b)] Si \struck{$\mathring{\ll}$} $X \mathrel{\mathring{\ll}} Y$, alors
    $\forall X' \in \widetilde{X}$, il existe $Y' \in \widetilde{Y}$ avec
    $X' \mathrel{\mathring{\ll}} Y'$ (NB $Y'$ unique par a)).
\end{itemize}

\emph{Mag 3'}, \emph{Mag 4'} comme \emph{Mag 3}, \emph{Mag 4}

\marginal{\uncertain{Ariadoshi} ;
relations entre $\leq$ et $\mathring{\ll}$ ;
Mag 3' relations entre $\mathring{\ll}$ et $|\circ|$ ;
Mag 4' relations entre $\leq$ et $|\circ|$}
\note{notes marginales écrites en oblique dans la marge gauche, en regard de Mag 1' -- Mag 4'.}

Ce n'est pas plus simple ni moins simple que Mag 1 -- Mag 4, mais la donnée
de $\mathring{\ll}$ paraît
\note{la phrase se poursuit à la page suivante.}

\page{30}
\note{p.~22 de l'auteur.}
plus « primaire » que celle de $\leq$ \add{et plus « intuitive »}, car

1°) Elle est « disjointe » de \struck{elle} donnée $\leq$, en fait on a
\[
  (X \leq Y \text{ et } X \mathrel{\mathring{\ll}} Y) \Longleftrightarrow X = Y ,
\]
\struck{et} et

2°) elle s'exprime plus naturellement dans Mag 3 = Mag 3'. On peut dire
\struck{\ill{}} \add{aussi} que les axiomes \struck{MAG} Mag 1' -- Mag 4' sont
à \uncertain{fonction} plus clairement séparées : Je vais les reformuler
légèrement. \add{(Notation intrinsèque des trois relations \uncertain{séparément})}

\emph{Mag 1'} $X \leq Y$ et $X \mathrel{\mathring{\ll}} Y$ sont des relations
d'ordre. $X \mathrel{|\circ|} Y$ est une relation symétrique. Antiréflexive.

\emph{Mag 2'} (Relations entre $\leq$ et $\mathring{\ll}$)
\begin{itemize}
  \item[a)] Si $X \mathrel{\mathring{\ll}} Y' \leq Y$, $Y'$ est uniquement
    \uncertain{déterminé} par $X, Y$
  \item[b)] Si $X \mathrel{\mathring{\ll}} Y$, alors $\forall X' \in \widetilde{X}$,
    $\exists Y' \in \widetilde{Y}$ tel que $X' \mathrel{\mathring{\ll}} Y'$.
\end{itemize}

\emph{Mag 3'} (Relations entre $\mathring{\ll}$ et $|\circ|$)
si $X \mathrel{|\circ|} Y$, $X' \mathrel{\mathring{\ll}} X$,
$Y' \mathrel{\mathring{\ll}} Y$, alors $X' \mathrel{|\circ|} Y'$

\emph{Mag 4'} (Relations entre $|\circ|$ et $\leq$) Si $X \leq Y$, $X \neq Y$,
alors \uncertain{non} $X \mathrel{|\circ|} Y$,

\marginal{Dans la version initiale, en somme, \uncertain{les} axiomes
\ill{} \ill{} \ill{} : \uncertain{Mag 1 : Mag 4} ; \uncertain{distingue} a) b) c)}
\note{note marginale oblique à gauche des axiomes, séparée d'eux par un trait vertical ; lecture très incertaine.}

Décidément, c'est la forme parfaite de définition !

Rappel des principales \struck{\ill{}} \add{\uncertain{phases}} \uncertain{parcourues}
depuis \uncertain{quinze} semaines, pour le choix des données de base,
depuis GF IV p.~52 (\struck{\ill{}} 13 juin)
\begin{itemize}
  \item[] \struck{\ill{}}, $\mathfrak{F}$, $\leq$, $\curlyeqprec$
  \item[] $\mathfrak{F}$, $\leq$, $\ll$
  \item[] $\mathcal{M}$, $\leq$, $\ll$, $\lessgtr$ (retour vers les origines)
  \item[] $\mathcal{M}$, $\leq$, $\ll$, $|\circ|$ \quad \add{(ou $\mathcal{M}$, $\leq$, $\ll$, $\|$)}
    Quand se rendre \uncertain{l'innovation} d'avant-hier faisant tomber tous les
    fouillis des axiomes de compatibilité -- disjonction
  \item[] $\mathcal{M}$, $\leq$, $\mathring{\ll}$, $|\circ|$ aujourd'hui,
    l'épanouissement intégral !
\end{itemize}
\note{le premier symbole de la première ligne est un $\leq$ surmonté d'un trait ondulé ; le mot raturé devant lui est illisible. L'ajout « ou $\mathcal{M}, \leq, \ll, \|$ » est relié par un trait à la ligne des $|\circ|$.}

\page{31}
\note{p.~23 de l'auteur.}
Ce choix \struck{d'axiomes} de données de base
\[
  \leq,\ \mathring{\ll},\ |\circ|
\]
fixe le magasin, pour ainsi dire, par le choix correspondant pour
\struck{$\mathfrak{F}$} un atelier $\mathfrak{F}$. Seule des trois relations
primitives, $\leq$ passe à $\mathfrak{F}$. Donc partant de $\mathfrak{F}$ tout
seul des lors \uncertain{pourrait} \uncertain{intégral}, il vaut mieux
prendre \struck{$(\mathcal{M}, \mathfrak{F}$}
\[
  (\mathcal{M}, \mathfrak{F} \mid \lhd, \mathring{\ll}, |\circ|)
\]
où $\lhd$ « relation d'incidence » entre $\mathcal{M}$ et $\mathfrak{F}$,
\add{et} $\mathring{\ll}$, $|\circ|$ relations dans $\mathcal{M}$ (la relation
$\leq$, dans $\mathcal{M}$ ou dans $\mathfrak{F}$ « sortant » de \struck{$\lhd$}
$\lhd$ \add{donnée}).

Voyons la forme que prennent les axiomes pour un « magasin strictement
divisible »* (cf.\ p.~7). La relation $|\circ|$ est ici redondante,
\struck{D'ailleurs} sur $\mathcal{M}$, on a la donnée (d'une relation $\|$)
sur $\mathcal{L}$. D'autre part, en termes de \struck{\ill{}} $\leq$,
$\mathring{\ll}$, on définit $\mathcal{L}$ comme
\[
  \mathcal{L} = \{ X \in \mathcal{M} \mid X \text{ minimal pour } \leq
  \text{ et pour } \mathring{\ll} \} .
\]
D'où données
\[
  (\mathcal{M}, \leq, \mathring{\ll}, \|) \qquad \text{avec}
\]

\emph{Mag div 1} $\leq$, $\mathring{\ll}$ sont des relations d'ordre sur
$\mathcal{M}$. La relation $\|$ est une relation symét.\ et antiréflexive sur
\add{l'ens.\ $\mathcal{L}$ des éléments de $\mathcal{L}$} minimaux pour $\leq$ et
$\mathring{\ll}$.

\emph{Mag div 2} Comme Mag 2' p.~\uncertain{9} précédente.

\emph{Mag div 3} Si $X \leq Y$, $X \neq Y$, $x \mathrel{\mathring{\ll}} X$,
$y \mathrel{\mathring{\ll}} Y$, $x, y \in \mathcal{L}$, alors $x \| y$.
\struck{\ill{}}
\note{l'ajout interlinéaire de Mag div 1 se lit mal ; la lecture « l'ens.\ $\mathcal{L}$ des éléments » est conjecturale. Le renvoi « p.~9 » est peut-être « p.~précédente » seul, la page précédente portant Mag 2'.}

\marginal{* Magasins où l'un des \uncertain{trois} \uncertain{derniers} \ill{}}
\marginal{On retrouve $X \mathrel{|\circ|} Y$ \ill{} : $X \mathrel{|\circ|} Y
\Leftrightarrow \forall x \mathrel{\mathring{\ll}} X$,
$y \mathrel{\mathring{\ll}} Y$, $x \| y$}
\note{deux notes marginales obliques, à gauche ; la seconde est reliée à Mag div 3.}

\page{32}
\note{p.~24 de l'auteur (le chiffre est récrit sur un autre).}
Cas \uncertain{quasi-ensembliste} où la relation $x \| y$ équivaut à $x \neq y$.
Données $(\mathcal{M}, \leq, \mathring{\ll})$, avec axiomes

\emph{Mag quens 1} $\leq$ et $\mathring{\ll}$ sont des relations d'ordre

\emph{Mag quens 2} Comme Mag 2' p.~21

\emph{Mag quens 3} $\forall X \in \mathcal{M}$, $\exists x \in \mathcal{L}$,
avec $x \mathrel{\mathring{\ll}} \mathcal{M}$.
\note{lire sans doute $x \mathrel{\mathring{\ll}} X$ ; la page porte bien un $M$ majuscule.}

\marginal{On retrouve $X \mathrel{|\circ|} Y$ par \ill{} :
$X \mathrel{|\circ|} Y \Leftrightarrow \mathrm{Omb}^{\circ}(X) \cap
\mathrm{Omb}^{\circ}(Y) = \emptyset$, où
$\mathrm{Omb}^{\circ}(Z) = \{ x \in \mathcal{L} \mid x \mathrel{\mathring{\ll}} Z \}$}
\note{note marginale oblique dans l'angle supérieur gauche ; la lecture des ensembles $\mathrm{Omb}^{\circ}$ est incertaine.}

Enfin, les magasins ensemblistes (les seuls \uncertain{ici} \uncertain{examinés}
auxquels \uncertain{jusque} aient travaillé \uncertain{vraiment} \uncertain{moi} !)
sont réduits à la structure de la p.~9, à laquelle je n'ai rien à ajouter,
sinon pour expliciter les structures $\leq$, $\ll$ comme induite par les
structures ensemblistes correspondantes dans $\mathrm{Fig}(\mathcal{L})$, et la
\struck{\ill{}} relation
\[
  X \mathrel{|\circ|} Y \Longleftrightarrow X^{\circ} \cap Y^{\circ} = \emptyset
\]
où $X^{\circ} = X \setminus \partial X$, $Y^{\circ} = Y \setminus \partial Y$,
constructions ensemblistes habituelles.

\begin{itemize}
  \item[] Mag 1 -- 4 (p.~22) Mag L 1, L 2, (p.~6.7) Mag at (p.~10) :
    pour des « \struck{\ill{}} magasins »
  \item[] \struck{Mag At 1 -- 5} (p.~11, 12) At 1 -- At 5 + Mag 1 -- \struck{4}
    \add{4} --- ateliers
  \item[] (p.~13, 14) At M 1 -- At M 5 + Mag 1 -- 4 --- mixtes
  \item[] (p.~23) Mag div 1 -- 3 (en lieu de Mag 1 -- 4, + L 2) --- magasins
  \item[] Mag quens 1 -- 3 (p.~24)
  \item[] Mag ens 1 -- 2 (p.~9)
  \item[] Variante polyidéale Mag 1 -- 4 + Mag pol (+ Mag at pol) (pages 15 -- 19) !
    Voir aussi Mag 1', Mag 2' p.~30, 31
\end{itemize}
\note{récapitulatif des systèmes d'axiomes, en bas de page ; les renvois aux pages sont ceux de l'auteur. Deux traits horizontaux à droite des trois dernières lignes restent sans libellé. Un long trait courbe relie « Mag at » à « Mag 1 -- 4, + L 2 ». Au lieu de « p.~30, 31 », la page porte peut-être « p.~31, 33 » récrit.}

\page{33}
\note{p.~25 de l'auteur.}
\emph{Magasins triviaux.} Déf : Un magasin est dit « trivial » si
$X \mathrel{\mathring{\ll}} Y \Leftrightarrow X = Y$, et si
$X \mathrel{|\circ|} Y \Leftrightarrow X \neq Y$, donc la seule donnée qui
\uncertain{reste} à déterminer est $X \leq Y$. Mais on vérifie \add{trivialement}
que pour \struck{tout} ens.\ ordonné $(\mathcal{M}, \leq)$, si on prend
\struck{$X \mathrel{\mathring{\ll}} Y$}
\[
  X \mathrel{\mathring{\ll}} Y \overset{\mathrm{def}}{\Longleftrightarrow} X = Y ,
  \qquad X \mathrel{|\circ|} Y \overset{\mathrm{def}}{\Longleftrightarrow} X \neq Y ,
\]
on trouve \add{\uncertain{que}} une structure de magasin.
\struck{Ce magasin ne satisfait Mag L 1 que si \ill{} $\mathcal{L} =$ ens.\ des
éléments minimaux}
\add{\uncertain{qui est} \uncertain{il}} \struck{\ill{}} \uncertain{On a} :
\[
  \begin{array}{lll}
    X \ll Y & \Longleftrightarrow & X \leq Y \\
    X \lessgtr Y & & \text{toujours satisfait} \\
    X \| Y & \text{ssi} & \widetilde{X} \cap \widetilde{Y} = \emptyset
      \text{ i.e. } \{X, Y\} \text{ pas minoré.}
  \end{array}
\]
$\mathcal{L}$ = ens.\ des éléments minimaux de $X$
\[
  \begin{array}{l}
    \mathrm{omb}(X) = \{ x \in \mathcal{L} \mid x \leq X \} \\
    \mathrm{omb}(X)^{\circ} = \text{\struck{$\{ x \in \mathcal{L} \mid x < X \}$}}
  \end{array}
  \left\{
  \begin{array}{ll}
    \emptyset & \text{si } X \notin \mathcal{L} \\
    \mathrm{omb}(X) = \{X\} & \text{si } X \in \mathcal{L}
  \end{array}
  \right.
\]
Donc la condition \emph{Mag L 1}, i.e.\ $\mathrm{omb}(X)^{\circ} \neq \emptyset$
$\forall X$, signifie que tout $X$ est minimal i.e.\ \uncertain{la} relation
d'ordre est discrète. On voit \uncertain{donc} qu'un magasin trivial ne satisfait
jamais Mag L 1, et \uncertain{même} \uncertain{guère} Mag L 2 --- \uncertain{sauf}
s'il est \uncertain{discret}-trivial ($X \leq Y \Leftrightarrow X = Y$).

\marginal{\emph{NB} Conditions \uncertain{équivalentes} pour un magasin
$\mathcal{M}$ : (i) \ill{} (ii) \ill{} (iii) $X \neq Y \Rightarrow$ \ill{}
$\mathrm{Fig}(\mathcal{M})$ \ill{} (iv) $\mathcal{M} \in \mathrm{Fig}$ \ill{}
(v) $X \neq Y \Rightarrow X \mathrel{|\circ|} Y$ \ill{} \ill{} ;
Ces conditions \ill{} \ill{} (vi) $X \mathrel{\mathring{\ll}} Y \Rightarrow X = Y$
(vii) $X \ll Y \Rightarrow X \leq Y$ \ill{} \ill{} magasins \ill{} \ill{}
\uncertain{triviaux} \ill{}}
\note{deux blocs de notes marginales très serrées, écrits en oblique dans la marge gauche sur la moitié inférieure de la page, et numérotés (i) à (vii) ; seuls quelques fragments s'en lisent, rendus ici sans garantie d'ordre.}

\emph{Figure indexée} par un ens.\ ordonné $I$ : figure $F$, avec \uncertain{donnée}
d'un ordonné
\[
  I \xrightarrow{\ \sim\ } \widetilde{F}, \qquad i \longmapsto X_{i}
\]
\note{l'argument se poursuit à la page suivante. Un symbole biffé précède $I$.}

\page{34}
\note{p.~26 de l'auteur.}
Ainsi, une figure indexée \add{par $I$}, définit une application
\[
  \varphi : I \longrightarrow \mathcal{M}
\]
qui est \emph{croissante} pour $\leq$ sur $\mathcal{M}$, \struck{Quelle est}
$F$ est déterminé par $\varphi$ \struck{\ill{}} \uncertain{avec}
\[
  \widetilde{F} = \varphi(I) .
\]
Quelles conditions sur $\varphi$, pour que ça \uncertain{corresponde} bien à une
figure indexée ? Il faut
\begin{itemize}
  \item[a)] $\varphi$ injectif, et \struck{une isom.} \add{induisant} \uncertain{un isom.}
    d'ens.\ ordonnés $I \simeq \varphi(I)$
\end{itemize}
et
\begin{itemize}
  \item[b)] $\varphi(I)$ est une figure. \struck{\ill{} dès que} \struck{\ill{}
    signifie deux choses}
\end{itemize}
\begin{itemize}
  \item[1°)] $\forall i, j \in I$, $X_{i} \lessgtr X_{j}$ \struck{i.e.} et
  \item[2°)] $\varphi(I)$ est une partie fermée de $\mathcal{M}$.
\end{itemize}
La première condition revient à celle-ci ; \struck{qui englobe aussi}
\begin{itemize}
  \item[b)] $\forall i \in I$, $\varphi \mid I_{\leq i} : I_{\leq i}
    \xrightarrow{\ \sim\ } \widetilde{X}_{i}$
  \item[c)] Si $i \neq j$, $X_{i} \mathrel{|\circ|} X_{j}$
\end{itemize}
\note{la condition c) est écrite telle quelle ; elle est corrigée dans la proposition qui suit.}

En d'autres termes, \uncertain{à} \uncertain{dire} $X \to X_{i}$,

\emph{Proposition.} Pour que $\varphi : I \to \mathcal{M}$, \add{définisse une
figure indexée par $I$,} il faut et il suffit que
\begin{itemize}
  \item[a)] $\varphi$ \struck{injectif} croissante (dans $\mathcal{M}$, $\leq$) et
    \emph{injection}.
  \item[b)] $\forall i \in I$, l'application induite
    \[
      \widetilde{i} = I_{\leq i} \xrightarrow{\ \sim\ } \widetilde{X}_{i} =
      \mathcal{M}_{\leq X_{i}}
    \]
    soit \struck{bijective} \add{bijection} \struck{(\ill{} \ill{})}
    \uncertain{(cette \ill{} ensembles ordonnés)} $=$ un isom.\ d'ensembles ordonnés)
  \item[c)] Si $i \neq j$ \add{(i.e.\ $i \mathrel{|\circ|} j$ pour la structure de
    magasin triviale sur $I$ !)} on a $X_{i} \mathrel{|\circ|} X_{j}$
\end{itemize}
\note{la proposition est marquée d'un trait vertical en marge ; dans b), la mention biffée sous la flèche n'est pas lisible.}

\page{35}
\note{p.~27 de l'auteur.}
En effet, b) implique déjà que $\varphi(I)$ est fermé. Il reste : prouver que
\struck{$i \neq j \Rightarrow$ \ill{} $X_{i}$} pour $i, j \in I$,
i.e.\ $X_{i} \lessgtr X_{j}$. \emph{Donc} que
\[
  X' \in \widetilde{X}_{i} \setminus \widetilde{X}_{i} \cap \widetilde{X}_{j},\
  Y' \in \widetilde{X}_{j} \setminus \widetilde{X}_{i} \cap \widetilde{X}_{j}
  \Longrightarrow X' \mathrel{|\circ|} Y' .
\]
Or par b) on a \ill{} $X' = X_{i'}$, avec $i' \leq i$, idem pour
$Y' = X_{j'}$, avec $j' \leq j$, d'ailleurs comme $X' \neq Y'$, on a
$i' \neq j'$ par \uncertain{a)}. \emph{Donc} on a $X_{i'} \mathrel{|\circ|} X_{j'}$
par c), qed.

\section{\struck{Morphismes de \ill{}} Variantes polyidéales}

On suppose maintenant l'ens.\ ordonné $I$ polyidéal. On veut se borner aux
figures polyidéales. \struck{\ill{}} Les \struck{deux \ill{}} n'y \uncertain{pas}
\uncertain{à} \uncertain{rien} faire ; \struck{\ill{}} \uncertain{à} \uncertain{vrai}
\uncertain{\ill{}tiquement}, puisque la condition \uncertain{polyidéale} sur une
figure se lit sur \uncertain{le} magasin $\widetilde{F}$, et signifie juste que
celle-ci est un ens.\ ordonné polyidéal.
\note{la première phrase de ce paragraphe se lit mal ; le symbole biffé devant $\widetilde{F}$ est illisible.}

\section{Morphismes de magasins $\varphi : \mathcal{M} \to \mathcal{M}'$ :}

application telle que si $X, Y \in \mathcal{M}$, $X' = \varphi(X)$,
$Y' = \varphi(Y)$, on ait
\begin{itemize}
  \item[a)] $X \leq Y \Longrightarrow X' \leq Y'$
  \item[b)] $X \mathrel{\mathring{\ll}} Y \Longrightarrow X' \mathrel{\mathring{\ll}} Y'$
  \item[c)] $X \mathrel{|\circ|} Y \Longrightarrow X' \mathrel{|\circ|} Y'$
    [+ (d) plus loin \uncertain{peut-être}.]
\end{itemize}
\marginal{$\varphi$ croissante pour $\leq$, $\mathring{\ll}$ (donc aussi pour
$\ll$) (Est-ce trop exigeant ??) Ça semble raisonnable dans le cas d'un
plongement de magasins seulement. On peut se demander si
$X \lessgtr Y \Longrightarrow X' \lessgtr Y'$, ce qui est moins exigeant…}
\note{remarque écrite à droite des conditions a) -- c) et au-dessous ; une variante biffée de « $\Longrightarrow$ » est récrite.}

\uncertain{on en} conclut
\[
  X \ll Y \Longrightarrow X' \ll Y'
\]
\uncertain{avec} \uncertain{aussi}
\[
  X \lessgtr Y \overset{?}{\Longrightarrow} X' \lessgtr Y' , \qquad
  X \| Y \overset{?}{\Longrightarrow} X' \lessgtr Y'
\]
\note{les deux dernières implications portent un point d'interrogation au-dessus de la flèche ; la seconde a bien $\lessgtr$ à droite, comme la première.}

\marginal{\uncertain{Peut-être} on \ill{} \ill{} \ill{} \ill{} \ill{} \ill{} par
\ill{} \ill{} \ill{} \uncertain{plus} \ill{} \ill{} ! Mais \ill{} \ill{} \ill{}
\ill{} \ill{} \ill{} \ill{}}
\note{note marginale oblique, en bas à gauche de la page, presque entièrement illisible.}

\marginal{\emph{NB} Si $\mathcal{M}$ est un magasin trivial, ces conditions
\uncertain{signifient} \uncertain{que} $\varphi$ \uncertain{injective}…}
\note{la dernière note, en pied de page, s'interrompt au bord de la feuille.}

\page{36}
\note{p.~28 de l'auteur.}
\struck{Il semble qu'il manque une condition pour \ill{} \ill{}}
\note{deux lignes encadrées et hachurées en tête de page.}
Il semble qu'on ait besoin aussi de la condition
\begin{itemize}
  \item[d)] $\forall X \in \mathcal{M}$, $\varphi$ induit une \emph{bijection}
    \[
      \widetilde{X} \xrightarrow{\ \sim\ } \widetilde{X}' \qquad \text{où }
      X' = \varphi(X) .
    \]
\end{itemize}

\marginal{\uncertain{11.7.86} Cette condition d) devrait \uncertain{peut-être}
\uncertain{être} \ill{} : \uncertain{en} \uncertain{fait} \ill{} \ill{} \ill{}
\ill{} \ill{} jamais \ill{}}
\note{note marginale oblique dans la marge gauche, en regard de d) ; la date se lit « 11.7.86 », postérieure à la fourchette du dossier, ou peut-être « 1.7.86 ».}

\emph{Corollaire 1.} Les conditions a), c), d) impliquent que
$\forall F \in \mathrm{Fig}(\mathcal{M})$, on a $\varphi(F) \in \mathrm{Fig}(\mathcal{M}')$,
et $\varphi \mid F$ \struck{est injective} \add{(\uncertain{indexé} par $I = \widetilde{F}$)}
et \uncertain{est} \uncertain{un} isom.\ d'ens.\ ordonnés. Immédiat par la
proposition précédente, si $\varphi$ est injective. \struck{Dans le cas général,}
Il suffit de vérifier que $\varphi \mid F$ est injective. Mais si $X, Y \in F$,
\struck{et} $X \neq Y$, on a $X \mathrel{|\circ|} Y$, donc
$\varphi(X) \mathrel{|\circ|} \varphi(Y)$ par c), donc
$\varphi(X) \neq \varphi(Y)$ OK ! \struck{M \ill{} que} (C'est la propr.\ qui
\uncertain{montre} en plus que $\varphi$ est un isom.\ ordonné…)

\emph{Corollaire 2.} \struck{Si $X \lessgtr Y$} Si $X' = \varphi(X)$,
$Y' = \varphi(Y)$
\[
  X \lessgtr Y \Longrightarrow X' \lessgtr Y' , \qquad
  X \| Y \Longrightarrow X' \| Y' .
\]
En effet, $X \lessgtr Y$ signifie que $X, Y$ font partie d'une figure $F$, et on
applique cor.~1. \uncertain{Plus} \struck{Si $X \| Y$, $X'$ \ill{} \ill{}
$\mathrm{Inf}(X, Y)$}

\uncertain{précisément}

\emph{Cor.~3.} Soient $F$ figure dans $\mathcal{M}$. Alors pour
$X, Y \in \widetilde{F}$, et $X' = \varphi(X)$, $Y' = \varphi(Y)$, on a
\[
  \begin{array}{l}
    \text{\struck{$X \lessgtr Y \Longleftrightarrow X' \lessgtr Y'$}} \\
    \text{\struck{$X \| Y$}} \\
    X \leq Y \Longleftrightarrow X' \leq Y' \\
    \text{\struck{$X \ill{}$}} \\
    X \| Y \Longleftrightarrow X' \| Y'
  \end{array}
\]
\marginal{\uncertain{La} généralisation \uncertain{aux} \uncertain{deux} figures
quelconques de $F$}
\note{le symbole biffé de l'avant-dernière ligne est noirci ; un astérisque suit le tableau.}

La première équivalence \struck{provient} déjà vue (cor.~1). La deuxième provient
du fait que $X \| Y$ \struck{$\{X, Y\}$} \uncertain{non} \ill{} \ill{} dans
l'ens.\ ordonné $I = F$ \uncertain{signifie} $\{X, Y\}$ \uncertain{non}
minoré) \uncertain{auquel} est isom.\ $I' = \varphi(I)$.
\note{les dernières lignes se lisent mal ; l'ordre des mots y est incertain.}

\page{37}
\note{p.~29 de l'auteur.}
\emph{Cor.~4.} Soient $F, G \in \mathrm{Fig}(\mathcal{M})$, avec $F \lessgtr G$.
Alors \uncertain{soient} $F' = \varphi(F)$, et $G' = \varphi(G)$, on a
\[
  F' \lessgtr G' , \quad \varphi(F \cup G) = \varphi(F) \cup \varphi(G) ,
  \quad \varphi(F \cap G) = \varphi(F) \cap \varphi(G)
\]
\emph{Dém.} On applique cor.~1 à $F \cup G$.

\emph{Cor.~5.} Si $F \ll G$, alors $\underset{F'}{\varphi(F)} \ll
\underset{G'}{\varphi(G)}$, et on a \uncertain{commutativité} de

\begin{tikzcd}
  \widetilde{F} \arrow[r] \arrow[d, "\varphi"'] & \widetilde{G} \arrow[d, "\varphi"] \\
  \widetilde{F}' \arrow[r] & \widetilde{G}'
\end{tikzcd}

\emph{Cor.~6.} Compatibilité de $\varphi$ avec l'opération des
\uncertain{affinements} induit sur les figures.

Je n'ai pas l'impression d'avoir bien saisi le rôle exact de l'hypothèse c), qui
semble très forte. Dans le cas d'une opération d'image inverse
\uncertain{\ill{}} par une application ensembliste, elle semble supposer
que cette dernière est \emph{surjective}. Mais dans ce cas, $\varphi$
\add{serait} \struck{est} \uncertain{injection}, ce qui n'est pas \uncertain{du}
\uncertain{tout} le \uncertain{cas} \uncertain{en} \uncertain{général},
\struck{\ill{} \ill{}}

\marginal{?? on \uncertain{devrait} \uncertain{plutôt} \uncertain{supposer}
\ill{} \ill{} que $\varphi(X) \mathrel{|\circ|} \varphi(Y) \Rightarrow
X \mathrel{|\circ|} Y$}
\note{note marginale oblique en bas à gauche ; une double flèche la relie au texte.}

\emph{Exemple.} 1) Considérons l'application
\[
  \mathcal{M} \longrightarrow \text{Figél}(\mathcal{M}) \qquad
  X \longmapsto \mathrm{Multombs}(X)
\]
\add{figures ensemblistes ici élémentaires}
C'est un \uncertain{plongement} \uncertain{compatible} pour $\leq$, $\ll$,
$\lessgtr$, qu'en est-il pour $|\circ|$ ?
On a
\[
  \begin{array}{l}
    |\mathrm{Multombs}(X)| = \mathrm{Omb}(X) = \mathcal{M}_{\ll X} \\
    |\mathrm{Multombs}(X)|^{\circ} = \mathrm{Omb}^{\circ}(X) =
      \mathcal{M}_{\mathring{\ll} X}
  \end{array}
\]
\note{les noms « Figél » et « Multombs » sont des lectures incertaines de mots abrégés (figures élémentaires, multi-ombres ?) ; le premier porte un « él » récrit, et l'ajout « figures ensemblistes ici élémentaires » est écrit au-dessous, en partie biffé.}

\page{38}
\note{p.~30 de l'auteur.}
donc
\[
  \text{\struck{$\mathrm{Multombs}(X) \mathrel{|\circ|} \mathrm{Multombs}(Y) \Longleftrightarrow$}}
\]
\[
  \begin{array}{rl}
    \mathrm{Multombs}(X) \mathrel{|\circ|} \mathrm{Multombs}(Y)
    & \Longleftrightarrow \mathrm{Omb}^{\circ}(X) \cap \mathrm{Omb}^{\circ}(Y) = \emptyset \\
    & \Longleftrightarrow \nexists Z \in \mathcal{M}, \text{ avec }
      Z \mathrel{\mathring{\ll}} X,\ Z \mathrel{\mathring{\ll}} Y .
  \end{array}
\]
\note{avant $\mathrm{Omb}^{\circ}(X)$, un premier essai $\mathcal{M}_{\mathring{\ll} X}$ est biffé ; à gauche de la seconde ligne, « i.e.\ \ill{} » est biffé.}

\struck{Est-ce} On a bien
\[
  X \mathrel{|\circ|} Y \Longrightarrow \mathrm{Multombs}(X) \mathrel{|\circ|}
  \mathrm{Multombs}(Y)
\]
i.e.\ $\{X, Y\}$ pas minoré dans $\mathcal{M}$, $\mathring{\ll}$. \uncertain{Car}
si on avait $Z \mathrel{\mathring{\ll}} X$, $Z \mathrel{\mathring{\ll}} Y$, on
aurait $Z \mathrel{|\circ|} Z$ par Mag 3 b), absurde par Mag 3 a).
\note{une ligne fine relie un $|\circ|$ écrit en marge au signe barré dans « $Z \mathrel{|\circ|} Z$ », écrit d'abord $Z \mathrel{\mathring{\ll}} Z$ ; le passage est encadré.}

Mais l'inverse est-il tjrs vrai i.e.\ $X \mathrel{|\circ|} Y$ peut-il s'exprimer en
termes de la relation d'ordre $\mathring{\ll}$ ?

\struck{Ça ne \uncertain{marche} \ill{} \ill{} C'est OK dans \ill{}}
Mais notons que si $x, y \in \mathcal{L}$, $x \neq y$, alors $\{x, y\}$ n'est pas
minoré dans $\mathcal{M}$ pour $\mathring{\ll}$, mais on n'aura pas pour autant,
en général, $x \mathrel{|\circ|} y$ i.e.\ $x \| y$ ! (Même dans le cas des
magasins divisibles.)

\emph{Donc} le fait que $|\circ|$ soit induit par $\mathring{\ll}$ s'\uncertain{exprime} :
une hypothèse de nature « quasi-ensembliste », mais sans \uncertain{référence}
particulière aux \uncertain{lieux}.

\page{39}
\note{p.~31 de l'auteur.}
On peut dès lors se demander quels axiomes \uncertain{mettre} sur
$(\mathcal{M}, \leq, \mathring{\ll})$, pour que la relation
\[
  (*) \qquad X \mathrel{|\circ|} Y \Longleftrightarrow \{X, Y\}
  \text{ non minoré dans } \mathcal{M}_{\mathring{\ll}}
\]
donne, avec $\leq$, $\mathring{\ll}$, un magasin. Les axiomes Mag 1, Mag 2
(p.~22), \uncertain{on} \uncertain{a} \uncertain{mis} \uncertain{avec} les ${}'$)
\uncertain{sont} \uncertain{clairs}, \struck{Mag 3' \ill{}} \uncertain{pour}
$\ll$, $\mathring{\ll}$, \uncertain{et} Mag 3' est automatique.
Mag 4' \struck{\uncertain{résulte} de Mag 2' a) ! \ill{} \ill{}}.
\note{la fin de la ligne biffée porte « Mag 2' a) ! », le 2 récrit sur un autre chiffre.}

\struck{Autrement}

\emph{Proposition.} Les magasins $(\mathcal{M}, \leq, \mathring{\ll})$ (avec
$|\circ|$ définissable par \struck{\ill{}} \uncertain{ci-dessus}) sont
caractérisés par les conditions\,:
\begin{itemize}
  \item[] \emph{Mag 1'} $\leq$, $\mathring{\ll}$ des relations d'ordre
  \item[] \emph{Mag 2'} Comme \uncertain{dans} \uncertain{le} \uncertain{magasin} p.~22.
\end{itemize}
\struck{Mag 3' \ill{} \ill{}, \ill{} $\mathring{\ll}$ \ill{}, $Y$}
\note{la proposition est marquée d'un trait vertical en marge ; la ligne Mag 3' qui la suivait est rayée de hachures.}

C'est \uncertain{pratiquement} aussi simple que le cas ensembliste (p.~9) ! Il
n'y a pas \uncertain{à} \uncertain{mettre} \uncertain{en} fait d'axiome
Mag quens, comme page 24. Finalement, on a une hiérarchie de structures de plus en
plus particulières, des notions \uncertain{de} magasins\,:
\[
  \begin{array}{lll}
    \text{Magasins généraux} & \leq,\ \mathring{\ll},\ |\circ| & \text{Mag 1 -- Mag 4} \\
    \text{Magasins} & \leq,\ \mathring{\ll} & \text{Mag 1', Mag 2' = Mag 2}
      \ (\text{p.~22}) \ (\text{\uncertain{cette} page}) \\
    \text{Magasins quasi-ensemblistes} & & \text{Mag 1', Mag 2', Mag 3'} \ (\text{p.~24}) \\
    \text{Magasins ensemblistes} & & \text{Magens 1, Magens 2} \ (\text{p.~9})
  \end{array}
\]
\marginal{Mettre un bon \uncertain{principe} ; \uncertain{et} \uncertain{c'est}
\uncertain{là} \uncertain{un} bon principe}
\note{les quatre lignes sont reliées par des doubles flèches verticales $\Uparrow$ ; à gauche, deux accolades portent la note marginale.}

On sera amené \uncertain{aussi} à regarder les magasins
\uncertain{d'autres} \uncertain{structures} \uncertain{(qui} interviennent
entre les magasins généraux et les ensemblistes…) \uncertain{qui} \uncertain{auront}
$\leq$, $|\circ|$, \ill{}
\note{la dernière phrase, serrée en pied de page et récrite entre les lignes, est d'ordre incertain.}

\page{40}
\note{p.~32 de l'auteur.}
\emph{Exemple 2.} Considérons l'application
\[
  \varphi : \mathcal{M} \longrightarrow \text{Figél}(\mathcal{L}) \qquad
  X \longmapsto \mathrm{multombs}(X)
\]
\struck{\ill{} \ill{} que \ill{}} (\uncertain{comme} \ill{} \uncertain{à} \uncertain{valeurs}
dans $\text{Figél}$) \uncertain{Ce} \uncertain{n'est} défini (que si on suppose
\uncertain{que}) Mag L 1, i.e.\ $\forall X \in \mathcal{M}$,
$\mathrm{omb}(X)^{\circ} \neq \emptyset$.

Dans ce cas, on a \uncertain{bien}
\[
  \begin{array}{l}
    \text{\struck{$X \mathrel{|\circ|} Y \Longrightarrow \mathrm{omb}(X)$}} \\
    |\mathrm{multombs}(X)| = \mathrm{omb}(X) \\
    |\mathrm{multombs}(X)|^{\circ} = \mathrm{omb}(X)^{\circ}
  \end{array}
\]
et
\[
  \mathrm{multombs}(X) \mathrel{|\circ|} \mathrm{multombs}(Y)
  \Longleftrightarrow \mathrm{omb}(X)^{\circ} \cap \mathrm{omb}(Y)^{\circ} = \emptyset
\]
Donc on a bien
\[
  X \mathrel{|\circ|} Y \Longrightarrow \mathrm{multombs}(X) \mathrel{|\circ|}
  \mathrm{multombs}(Y) .
\]

\struck{\ill{} \ill{} \ill{}}
Ces deux exemples nous montrent un \uncertain{peu} \uncertain{mieux} que la notion de
\uncertain{comorphismes} \uncertain{donnée} \uncertain{à} \uncertain{la} p.~27,
n'est pas trop idiote (\uncertain{même} s'il convient \uncertain{sans} \uncertain{doute}
de la \uncertain{généraliser}). Troisième exemple\,:

\emph{Exemple 3.} $f : E' \to E$ application ensembliste \uncertain{surjective}.
Elle induit un comorphisme \uncertain{injectif} des magasins ensemblistes associés
\[
  \text{Figél}(E) \longrightarrow \text{Figél}(E')
\]
\note{l'exemple 3 est marqué d'un trait vertical en marge ; $E'$ est récrit sur une autre lettre, et « surjective », « injectif » sont des lectures incertaines. L'argument se poursuit au-delà de la fin du lot.}

\end{document}
